% Part: incompleteness
% Chapter: representability-in-q
% Section: c

\documentclass[../../../include/open-logic-section]{subfiles}

\begin{document}

\olfileid{inc}{req}{cee}
\olsection{فنکشاں دا سیٹ $C$}

منّ لو $C$ فنکشاں دا سبھ توں نِکّا سیٹ اے جیہدے وچ ایہ شامل نیں:
\begin{enumerate}
\item $0$،
\item جانشین فنکشن،
\item جمع،
\item ضرب،
\item پروجیکشن فنکشن، تے
\item برابری دا امتیازی فنکشن، $\Char{=}$۔
\end{enumerate}
نالے ایہ سیٹ ایہناں عملّاں ہیٹھ بند اے:
\begin{enumerate}
\item فنکشاں دی ترکیب، تے
\item باقاعدہ فنکشاں اُتے غیر محدود تلاش۔
\end{enumerate}
چیتے رکھو، آخری پابندی دا مطلب بس ایہ اے کہ تسی $\mu$ عمل اوہدوں ای ورت سکدے او جدوں نتیجہ کُل فنکشن ہووے۔ ایہدی \emph{عام بازگشتی} فنکشاں دی تعریف نال ٹاکرا کرو۔ ایتھے جمع، ضرب تے $\Char{=}$ شامل کیتے نیں، تے بنیادی بازگشت چھڈ دِتّی اے۔

صاف اے کہ $C$ وچلا ہر فنکشن بازگشتی اے، کیوں جے جمع، ضرب تے $\Char{=}$ بازگشتی نیں۔ اسیں وکھاواں گے کہ اُلٹی گل وی سچ اے۔ ایہدا مطلب اے کہ $C$ وچلیاں باقی شےواں نال اسیں بنیادی بازگشت کر سکدے آں۔

\end{document}
